% !TeX root = ../../../金融经济学_王江_讲义及习题.tex
% 来源：教材图12.1；布局与配色参考431计算讲义有效前沿图。
% 用解析双曲线绘制，红色上支有效，灰色下支无效。
\begin{tikzpicture}[>=Stealth,x=1.65cm,y=1.1cm,font=\small]
 \fill[blue!5] plot[domain=-2:2.3,samples=85,variable=\t]
   ({sqrt(1+\t*\t)},{1.8+.7*\t})--(3.5,3.41)--(3.5,.4)--cycle;
 \draw[->] (0,0)--(4,0) node[right] {$\sigma$（风险）};
 \draw[->] (0,0)--(0,3.9) node[above] {$\bar r$（预期收益）};
 % 示意双曲线的两条渐近线为 y=1.8±0.7x。
 \draw[gray!75,dashed] (2.6,{1.8+.7*2.6})--(0,1.8)--(2.5,{1.8-.7*2.5});
 \draw[gray!65,thick,domain=-2:0,samples=60,variable=\t]
   plot ({sqrt(1+\t*\t)},{1.8+.7*\t});
 \draw[red!70!black,very thick,domain=0:2.3,samples=60,variable=\t]
   plot ({sqrt(1+\t*\t)},{1.8+.7*\t});
 \draw[densely dotted,gray] (0,1.8) node[left] {$a/c$}--(1,1.8)--(1,0) node[below] {$1/\sqrt c$};
 \fill (1,1.8) circle(2pt) node[above left] {MVP};
 % E 与 F 同均值，横向箭头指向该均值下的最小方差组合。
 \coordinate (efficient) at (1.72047,2.78);
 \coordinate (sameMean) at (3.16,2.78);
 \draw[dotted,gray] (1.72047,.82)--(1.72047,2.78);
 \draw[->,black!60,shorten <=4pt,shorten >=4pt] (sameMean)--(efficient);
 \fill[blue!60!black] (sameMean) circle(1.8pt) node[above right] {$F$};
 \fill[red!70!black] (efficient) circle(1.8pt) node[above left] {$E$};
 \fill[gray] (1.72047,.82) circle(1.8pt) node[below left] {$D$};
 \draw[->,black!60] (1.85,1.05)--(1.85,2.55);
 \node[red!70!black,above] at (2.5,3.4) {有效前沿 MVE};
 \node[gray] at (2.6,.55) {无效下半支};
 \node[blue!60!black,align=center] at (2.9,1.85) {可行组合\\可能所在区域};
\end{tikzpicture}
